% Section 9
\section{Open questions}\label{sec:open}

The proved results are stated in Sections~\ref{sec:rxc}--\ref{sec:inference}.
This section collects what is not proved: the conjectures and algorithmic
procedures that appear in the paper, and the open problems.

\subsection{Conjectures and algorithmic procedures}

\begin{itemize}
\item \textbf{CLR as the default interval} (Conjecture~\ref{conj:clr}). The
package's default confidence interval collects candidate endpoints by simulation
and computes its selection threshold over the collected set. It is an algorithm.
The reduction to the cited full-candidate CLR theorem is proved
(Lemma~\ref{lem:reduction}); the recovery of rarely optimal vertices, the
anti-conservative effect of computing the threshold over a subset, a growing
number of candidates, and uniform validity are open. Its coverage was at least
0.994 wherever Slater's condition held, and nominal at a point-identified design,
but simulation cannot make it a theorem. The variance-adaptive allowance computed
from a CLR interval, and the CLR branch union for contextual reinforcement,
inherit this status. Proving the conjecture would make the default a theorem.
\item \textbf{The sharp N2 set} (Conjecture~\ref{conj:N2}). Sufficiency of the
selection-expectation characterisation on a finite support is shown; necessity in
general is not.
\item \textbf{The balance point} (Conjecture~\ref{conj:C1}). That $x^*$ lies in the
composition hull in $2\times2$, which would complete the chain from signed slopes to
the within sign.
\item \textbf{Shared and district-specific effects in regime R3}: the asymptotic
results with a high-dimensional Gaussian approximation, under a rate condition on
the number of districts relative to their size.
\item \textbf{The wild-bootstrap calibration} of the pooling and curvature tests
is supported by simulation; its validity for a growing number of restrictions is
not proved.
\item \textbf{Spanning trees of pairwise beliefs}: that every spanning tree, not
only a star, reaches all directions of $K$ (checked for chains), and a related
conjecture that single-index contextual effects have dimension $C-1$.
\end{itemize}

\subsection{Open problems}

\begin{itemize}
\item \textbf{Sign rows under district rungs.} Sharpness and inference for the sign
and reinforcement rows under M1, M1($\delta$), M2 and L4. The rows are valid linear
rows, and the bound with them is an outer set, but neither its sharpness for the
regression nor the coverage of its intervals is proved. Until then the grouped
sign module is experimental and its intervals have no coverage theorem.
\item \textbf{Integer unit tables.} The integer--continuous gap beyond the
$2\times2$ case with margins only (Proposition~\ref{prop:integer}): with unit-level
restrictions, several columns or three or more groups.
\item \textbf{Inference on width gains.} Whether a belief tightens a target is
decided exactly at the population (Proposition~\ref{prop:R2}) and estimated at the
point estimate in samples. Inference on the width change needs joint treatment of
the base and augmented endpoints and of any estimated threshold; an interval for
either endpoint alone is not enough.
\item \textbf{Informative finite-sample inference with many small groups.}
Theorem~\ref{thm:finite} is valid in regime R2 but uninformative. Candidates are a
direct bound on the endpoint rather than on every constraint, variance-adaptive
(empirical-Bernstein) radii, a Bhatia--Davis version of the residual radius, and
uneven splits of the error level. A hierarchical model for district shifts would
change the model.
\item \textbf{Uniform limit theory under changing active sets.}
Proposition~\ref{prop:triangular} covers convergent localised slacks; a theory
uniform in the design sequence, when candidates stay within $O(n^{-1/2})$ of binding
without converging, is open. So are uniform validity near point identification for
the CLR interval (a Stoye-type adaptation) and the rate and limit of the endpoint
estimator when Slater's condition fails.
\item \textbf{The hull error.} Its rate, and inference for the population-hull
endpoint.
\item \textbf{Other targets and assumptions.} Inference for the superpopulation
target $B^*$; dependence within districts beyond the district effect; a joint
relaxation of the closed-population assumption for turnover; coverage of the
plug-in under misspecification; the characterisation of strict gains from joint
fitting and from shared district effects when $R\ge3$; refutability of sign,
ordering and local beliefs when $R\ge3$; an outer N2 LP with an estimated
regression; and a consistent series bound that would tighten $\ThO$ towards $\ThX$,
which conditions on $X$ only and so would not be sharp when size is informative.
\end{itemize}

\subsection{Future work}

Taking election periods as districts turns the shared-effects rung into the
assumption that contextual effects are unchanged between elections, which opens a
route to bounds on changes in ecological rates and to difference-in-differences
designs; that extension is left for future work.
